\uuid{NUGa}
\titre{Primitives usuelles}
\auteur{hanani ; mzari}
\organisation{exo7}
\licence{CC BY-NC-SA 4.0}
\datecreate{2020-01-20}
\niveau{L1}
\module{Analyse}
\chapitre{Calcul d'intégrales}
\sousChapitre{Primitives diverses}
\difficulte{1}
\theme{Primitives, primitives usuelles}
\exo7qcmid{lille-2.primitives.006}

\contenu{
\begin{question}
\qclassification{127.11}
\qtags{motcle=primitives usuelles}
Parmi les affirmations suivantes, cocher celles qui sont vraies. Sur $]0,+\infty[$, on a :
\begin{answers}
\bad{$\displaystyle \int \left(\frac{2}{x}+\mathrm{e}^{2x}\right)\mathrm{d}x=\frac{-2}{x^2}+2\mathrm{e}^{2x}+k$, où $k\in \Rr$.}
\good{$\displaystyle \int \left(\frac{2}{x}+\mathrm{e}^{2x}\right)\mathrm{d}x=\ln (x^2)+\frac{1}{2}\mathrm{e}^{2x}+k$, où $k\in \Rr$.}
\good{$\displaystyle \int \left(\frac{1}{2\sqrt{x}}-\sin x\right)\mathrm{d}x=\sqrt{x}+\cos x+k$, où $k\in \Rr$.}
\bad{$\displaystyle \int \left(\frac{1}{2\sqrt{x}}-\sin x\right)\mathrm{d}x=\sqrt{x}-\cos x+k$, où $k\in \Rr$.}
\end{answers}
\vskip3mm
\begin{explanations}
On a : $\displaystyle \int \frac{2\, \mathrm{d}x}{x}=\ln (x^2)+k$ et $\displaystyle \int\mathrm{e}^{2x}\mathrm{d}x=\frac{1}{2}\mathrm{e}^{2x}+k$. D'où, par linéarité,
$$\int \left(\frac{2}{x}+\mathrm{e}^{2x}\right)\mathrm{d}x=\ln (x^2)+\frac{1}{2}\mathrm{e}^{2x}+k.$$
De même, on vérifie que : $\displaystyle \int \left(\frac{1}{2\sqrt{x}}-\sin x\right)\mathrm{d}x=\sqrt{x}+\cos x+k$, où $k\in \Rr$.
\end{explanations}
\end{question}
}
